% abstract.tex, included by main.tex inside \begin{abstract}...\end{abstract}
% REVISION (advisor feedback, 2026-07): plain-language opening, the conceptual
% contribution stated first, claim scope corrected so that extraction is the measured
% outcome and sensitive-data leakage is an acknowledged assumption rather than a result.
Large language models are ranked and certified as safe using public benchmarks, and that
practice is only valid if the benchmark was absent from pre-training. When a benchmark
does appear in the training corpus, a condition known as \emph{benchmark contamination},
the usual concern is that scores overstate capability. A growing line of argument treats
contamination instead as a privacy warning sign, on the reasoning that contamination
reveals \emph{memorization}, and that memorization is the mechanism by which sensitive
training content escapes. That reasoning is rarely stated precisely, and contamination and
sensitive-data leakage are distinct phenomena, so we begin by decomposing it into four
explicit assumptions and identifying which are empirically testable. Our experiments
target the first and most basic of them: that a contamination score reflects retention of
a \emph{specific} item rather than how intrinsically predictable that item is. We test it
on Pythia models trained on the public Pile, so that training-set membership is ground
truth rather than an inferred label, and we compare four existing reference-free
detectors (LOSS, Min-K\%, Min-K\%++, and a compression ratio) against a per-item
extraction outcome. We propose no new detector, attack, or metric. Our contribution is a
controlled measurement of \emph{incremental} value: a pre-registered partial-correlation
and mediation analysis that asks what each detector predicts once the model's per-item
loss is held fixed. The three adjusted detectors add no positive predictive value beyond
loss. They are near-collinear with it (Spearman $0.74$ to $0.90$), and their residual
associations are null or, for the most loss-collinear detector, weakly negative in a
manner consistent with a statistical suppression artifact rather than genuine inverse
prediction. The absence of positive residual survives non-linear controls for loss and
survives deduplication, and is not explained by token frequency or by the
mostly-zero outcome. We read this as a measurement-validity result: the detectors an
auditor can afford carry little information about \emph{which} items leak that loss does
not already carry, which bounds what detector-based contamination auditing can support.
We do not demonstrate leakage of sensitive information. A PII measurement on the Enron
Emails subset of the Pile returned no detected leakage at our scale, and we report that
null rather than treating the sensitive-data link as established. \textbf{These results
are preliminary, obtained on the smallest ($160$M) Pythia model on CPU, which is the
regime least likely to exhibit memorization. The pipeline is built so that replication at
larger scale is a single configuration change.} All analyses are pre-registered and every
number is reproducible from a seeded script.
